% Unit: Computer Science A --- logging and testing (MC; \input after \texttt{cs-a/lambda.tex} in \texttt{main.tex}).

\runningheader{Computer Science A}{Unit 10: Testing and logging}{Page \thepage\ of \numpages}

\begingradingrange{csau10}

\question[4]
Henry came up with a few options for logging and presented it to Eric the SRE team lead. Which of the following option would be of highest priority to log?
\begin{choices}
  \choice 200 status codes
  \choice When the application is on low usage
  \CorrectChoice When the end user experiences a slow-down of the application
  \choice None of the above
\end{choices}
\begin{solution}
  User-visible latency and errors usually trump routine success metrics for incident response.
\end{solution}

\question[4]
When you write a test case, you should test every single method of an application.
\begin{choices}
  \CorrectChoice False
  \choice True
\end{choices}
\begin{solution}
  Prefer focused tests: one behavior or rule per test; use several \texttt{@Test} methods rather than one mega-test per production method.
\end{solution}

\question[4]
\texttt{@BeforeEach} will run every time, before each unit test.
\begin{choices}
  \CorrectChoice True
  \choice False
\end{choices}
\begin{solution}
  JUnit~5 \texttt{@BeforeEach} (JUnit~4 \texttt{@Before}) runs before each test method in that class.
\end{solution}

\question[4]
What is the package for JUnit?
\begin{choices}
  \choice \texttt{java.junit}
  \CorrectChoice \texttt{org.junit}
  \choice \texttt{javax.junit}
  \choice \texttt{org.oracle.java.junit}
  \choice \texttt{java.test.unit}
\end{choices}
\begin{solution}
  JUnit~4 lives under \texttt{org.junit}; JUnit~5 APIs are primarily \texttt{org.junit.jupiter}.
\end{solution}

\newpage

\question[4]
Predict the output of the following test:
\begin{lstlisting}[basicstyle=\ttfamily\footnotesize,frame=single,breaklines=true,aboveskip=4pt,belowskip=4pt]
public class ExpenseApprover {
  public boolean approveOrDenyExpense(Expense e) {
    if (e.amount < 100 && e.hasReceipt) {
      return true;
    } else {
      return false;
    }
  }
}

public class Test {
  ExpenseApprover ea = new ExpenseApprover();

  @Test
  public void testExpense() {
    Expense exp = new Expense();
    exp.setAmount(105);
    exp.setReceipt(true);
    Assert.assertFalse(ea.approveOrDenyExpense(exp));
  }
}
\end{lstlisting}
\begin{choices}
  \CorrectChoice Pass
  \choice Fail
  \choice Compilation error
  \choice Test doesn't run
\end{choices}
\begin{solution}
  Amount 105 fails \texttt{amount < 100}, so the method returns \texttt{false}; assertion expects \texttt{false}---passes.
\end{solution}

\endgradingrange{csau10}
